% GENERATED FILE. Edit canonical lessons, then run npm run generate:books.
% canonical-source-hash: fnv1a64:88d1e3c9a7dc0a2c
% canonical-lessons: LA-C455-conspicor, LA-C455-accedo, LA-C455-revertor, LA-C455-moror, LA-C455-commoror

\chapter{To catch sight of, To come near, To turn back, To delay, To stop for a while}
\label{ch:la-to-catch-sight-of-to-come-near-to-turn-back-to-delay-to-stop-for-a-while}
\hlchaptermodality{455}

\begin{chapteropening}
\textbf{By the end of this chapter:} \emph{I can talk about the end of a journey: catching sight of the town, coming near, turning back, delaying and stopping for a while.}

Complete the last lesson of chapter 455: I can talk about the end of a journey: catching sight of the town, coming near, turning back, delaying and stopping for a while.
\end{chapteropening}

\section[cōnspicor, cōnspicārī]{cōnspicor, cōnspicārī — to catch sight of, to see}

\label{lesson:LA-C455-conspicor}



\begin{quote}
Before the new one: say the Latin for to go round, then the Latin for a detour.
\end{quote}

\subsection*{You'll want to know: cōnspicor, cōnspicārī}
\textbf{cōnspicor, cōnspicārī} — \textquotedblleft{}to catch sight of, to see\textquotedblright{}.

\begin{quote}
Tandem viātōrēs mūrōs urbis cōnspicantur.
\end{quote}

\begin{quote}
\emph{At last the travellers catch sight of the city walls.}
\end{quote}

Deponent: \emph{cōnspicor, cōnspicārī, cōnspicātus sum}, with an accusative. A related active verb is \emph{cōnspiciō}.

\begin{cousinweb}
\emph{cōnspicor} is built from \emph{con-} and the root of \emph{speciō} 'to look' (Lewis \& Short), the root seen in \emph{spectō} 'to watch'. EtymDB lists it but gives no parent.
\end{cousinweb}

\begin{grammarlens}[title={What you've built}]
One more word for telling, explaining and describing at B1.
\end{grammarlens}

\subsection*{Your turn}
\begin{itemize}
\raggedright
  \item \emph{Say it:} \emph{cōnspicor, cōnspicārī}
  \item \emph{Say it:} \emph{cōnspicor, cōnspicārī}, once more
  \item \emph{Say it:} a sentence of your own with \emph{cōnspicor, cōnspicārī}
  \item \emph{Recall:} say the Latin for to go round, then the Latin for a detour, then say \emph{cōnspicor, cōnspicārī} again
\end{itemize}

\begin{tcolorbox}[breakable,colback=teal!4,colframe=teal!35!black,title={Before you move on}]
What does cōnspicor, cōnspicārī mean? (\textquotedblleft{}To catch sight of, to see\textquotedblright{}.) Say it once more.
\end{tcolorbox}
\section[accēdō, accēdere]{accēdō, accēdere — to come near, to approach; to be added}

\label{lesson:LA-C455-accedo}



\begin{quote}
Before the new one: say the Latin for a detour, then the Latin for to catch sight of.
\end{quote}

\subsection*{You'll want to know: accēdō, accēdere}
\textbf{accēdō, accēdere} — \textquotedblleft{}to come near, to approach; to be added\textquotedblright{}.

\begin{quote}
Nāvis ad portum accēdit.
\end{quote}

\begin{quote}
\emph{The ship comes near the harbour.}
\end{quote}

\emph{accēdō, accēdere, accessī, accessum}, with \emph{ad} and the accusative. English \emph{access} and \emph{accede} come from it.

\begin{cousinweb}
\emph{accēdō} is built from \emph{ad-} 'towards' and \emph{cēdō} 'to go' (EtymDB; Lewis \& Short). English \emph{access} and \emph{accede} come from it.
\end{cousinweb}

\begin{grammarlens}[title={What you've built}]
One more word for telling, explaining and describing at B1.
\end{grammarlens}

\subsection*{Your turn}
\begin{itemize}
\raggedright
  \item \emph{Say it:} \emph{accēdō, accēdere}
  \item \emph{Say it:} \emph{accēdō, accēdere}, once more
  \item \emph{Say it:} a sentence of your own with \emph{accēdō, accēdere}
  \item \emph{Recall:} say the Latin for a detour, then the Latin for to catch sight of, then say \emph{accēdō, accēdere} again
\end{itemize}

\begin{tcolorbox}[breakable,colback=teal!4,colframe=teal!35!black,title={Before you move on}]
What does accēdō, accēdere mean? (\textquotedblleft{}To come near, to approach; to be added\textquotedblright{}.) Say it once more.
\end{tcolorbox}
\section[revertor, revertī]{revertor, revertī — to turn back, to return}

\label{lesson:LA-C455-revertor}



\begin{quote}
Before the new one: say the Latin for to catch sight of, then the Latin for to come near.
\end{quote}

\subsection*{You'll want to know: revertor, revertī}
\textbf{revertor, revertī} — \textquotedblleft{}to turn back, to return\textquotedblright{}.

\begin{quote}
Via clausa est; itaque viātōrēs domum revertuntur.
\end{quote}

\begin{quote}
\emph{The road is closed; so the travellers turn back home.}
\end{quote}

Deponent in the present: \emph{revertor, revertī}, but the perfect is active: \emph{revertī}, \emph{reversus} as a participle. English \emph{revert} comes from it.

\begin{cousinweb}
\emph{revertor} is built from \emph{re-} 'back' and \emph{vertō} 'to turn' (EtymDB; Lewis \& Short, under \emph{revertō}). English \emph{revert} comes from it.
\end{cousinweb}

\begin{grammarlens}[title={What you've built}]
One more word for telling, explaining and describing at B1.
\end{grammarlens}

\subsection*{Your turn}
\begin{itemize}
\raggedright
  \item \emph{Say it:} \emph{revertor, revertī}
  \item \emph{Say it:} \emph{revertor, revertī}, once more
  \item \emph{Say it:} a sentence of your own with \emph{revertor, revertī}
  \item \emph{Recall:} say the Latin for to catch sight of, then the Latin for to come near, then say \emph{revertor, revertī} again
\end{itemize}

\begin{tcolorbox}[breakable,colback=teal!4,colframe=teal!35!black,title={Before you move on}]
What does revertor, revertī mean? (\textquotedblleft{}To turn back, to return\textquotedblright{}.) Say it once more.
\end{tcolorbox}
\section[moror, morārī]{moror, morārī — to delay, to linger; to hold up}

\label{lesson:LA-C455-moror}



\begin{quote}
Before the new one: say the Latin for to come near, then the Latin for to turn back.
\end{quote}

\subsection*{You'll want to know: moror, morārī}
\textbf{moror, morārī} — \textquotedblleft{}to delay, to linger; to hold up\textquotedblright{}.

\begin{quote}
Viātor in caupōnā nōn morātur, quod nox venit.
\end{quote}

\begin{quote}
\emph{The traveller does not linger at the inn, because night is coming.}
\end{quote}

Deponent: \emph{moror, morārī, morātus sum}, from \emph{mora} 'delay'. \emph{Nōlī morārī}, don't hang about. English \emph{moratorium} comes from the same root.

\begin{cousinweb}
\emph{moror} is built on \emph{mora} 'delay' (EtymDB; Lewis \& Short), a regular first-conjugation deponent made from the noun.
\end{cousinweb}

\begin{grammarlens}[title={What you've built}]
One more word for telling, explaining and describing at B1.
\end{grammarlens}

\subsection*{Your turn}
\begin{itemize}
\raggedright
  \item \emph{Say it:} \emph{moror, morārī}
  \item \emph{Say it:} \emph{moror, morārī}, once more
  \item \emph{Say it:} a sentence of your own with \emph{moror, morārī}
  \item \emph{Recall:} say the Latin for to come near, then the Latin for to turn back, then say \emph{moror, morārī} again
\end{itemize}

\begin{tcolorbox}[breakable,colback=teal!4,colframe=teal!35!black,title={Before you move on}]
What does moror, morārī mean? (\textquotedblleft{}To delay, to linger; to hold up\textquotedblright{}.) Say it once more.
\end{tcolorbox}
\section[commoror, commorārī]{commoror, commorārī — to stop for a while, to stay}

\label{lesson:LA-C455-commoror}



\begin{quote}
Before the new one: say the Latin for to turn back, then the Latin for to delay.
\end{quote}

\subsection*{You'll want to know: commoror, commorārī}
\textbf{commoror, commorārī} — \textquotedblleft{}to stop for a while, to stay\textquotedblright{}.

\begin{quote}
Mercātor trēs diēs in oppidō commorātur.
\end{quote}

\begin{quote}
\emph{The merchant stays three days in the town.}
\end{quote}

Deponent: \emph{commoror, commorārī, commorātus sum}. A short stay on the way, not an end.

\begin{cousinweb}
\emph{commoror} is built from \emph{con-} and \emph{moror} 'to delay, to linger' (EtymDB; Lewis \& Short).
\end{cousinweb}

\begin{grammarlens}[title={What you've built}]
One more word for telling, explaining and describing at B1.
\end{grammarlens}

\subsection*{Your turn}
\begin{itemize}
\raggedright
  \item \emph{Say it:} \emph{commoror, commorārī}
  \item \emph{Say it:} \emph{commoror, commorārī}, once more
  \item \emph{Say it:} a sentence of your own with \emph{commoror, commorārī}
  \item \emph{Recall:} say the Latin for to turn back, then the Latin for to delay, then say \emph{commoror, commorārī} again
\end{itemize}

\begin{tcolorbox}[breakable,colback=teal!4,colframe=teal!35!black,title={Before you move on}]
What does commoror, commorārī mean? (\textquotedblleft{}To stop for a while, to stay\textquotedblright{}.) Say it once more.
\end{tcolorbox}
